\chapter{黑箱之困——大语言模型的根本性局限与学术前沿}
\label{ch:limitations}
%% ============================================================

在前面的章节中，
我们见证了大语言模型从预测下一个词的简单目标出发，
一路攀升至超越人类博士的基准表现。
然而，技术进步的耀眼光芒不应遮蔽一个根本性的事实：
\textbf{我们仍然不真正理解这些系统为什么工作，
也不完全清楚它们在什么条件下会失败。}

本章的目标不是泼冷水，而是提供一幅诚实的技术图景——
LLM 在哪些维度上已经达到甚至超越了人类专家水平？
在哪些维度上存在\textbf{架构性的}、
而非仅凭规模堆叠就能解决的根本局限？
学术界在最近一年中
提出了哪些可能从根本上改变游戏规则的新思路？
以及，Yann LeCun 所倡导的世界模型路线
是否指向了一条超越当前范式的替代道路？

\section{悖论：超越博士的智能，无法解释的黑箱}
\label{sec:paradox}

2026 年初的 AI 能力画卷令人目眩。
让我们用冷静的数据审视这幅画卷的两面。

\subsection{令人震惊的一面}

\textbf{科学推理}。
GPQA Diamond 是一组由各领域博士专家编写的研究生级科学问题，
人类博士专家的准确率为 65\%，
而 Gemini 3.1 Pro 达到 94.3\%，
GPT-5.4 达到 92.0\%，
Claude Opus 4.6 在 mid-to-high 80s 范围~\cite{rein2024gpqa}。
\textbf{AI 在博士级科学推理上超越人类专家 25--30 个百分点}。

\textbf{数学竞赛}。
AIME 2025（美国数学邀请赛）上，
GPT-5.2 达到 \textbf{100\%}——
这个曾经被认为需要天才级数学直觉的基准
已经被前沿模型\textbf{完全饱和}。
Gemini Deep Think 在 2025 年国际数学奥林匹克（IMO）
中获得了正式金牌。

\textbf{编程竞赛}。
Gemini 3 Deep Think 在 Codeforces 上
达到 Elo 3455——
\textbf{地球上目前仅有 7 个人类的评分更高}。
从 o3 的 2724（2025 年 2 月）
到 Gemini 3 Deep Think 的 3455（2026 年 2 月），
一年之内进步了 700 多分。

\textbf{软件工程}。
SWE-bench Verified 上，
Claude Opus 4.6 达到 80.8\%。
从 2024 年 6 月 Sonnet 3.5 的 49.0\%
到 2026 年初的 80\%+，
20 个月内提升了 30 多个百分点。
在 METR Time Horizon 基准上，
Claude Opus 4.6 的 50\% 成功率时间地平线
达到 \textbf{14.5 小时}——
即它能完成相当于人类专家需要 14.5 小时才能完成的
软件工程任务。

\textbf{计算机操作}。
GPT-5.4 在 OSWorld（桌面自动化基准）上达到 75\%，
首次\textbf{超越人类表现}（人类 72.4\%）。

\textbf{职业知识}。
GPT-5.4 在 GDPval 上达到 83\%，
在 44 个职业领域
（会计、销售、工程等）
匹配或超越了人类专业人士。

\subsection{令人不安的一面}

然而，另一组数据讲述了截然不同的故事。

\textbf{真实职业任务}。
APEX-Agents 等面向真实职业任务的基准（2026 年初）
评估了 AI 在投资银行、管理咨询和公司法律中的实际工作任务：
\textbf{AI 模型在首次尝试中失败 76--82\%}，
即使给予 8 次重试机会，
也只能达到约 40\% 的成功率。

\textbf{人类最后的考试}。
``Humanity's Last Exam''（HLE）~\cite{hle2026nature}
由 2500 道跨数十个学科的题目组成，
2026 年 1 月发表于 \textit{Nature}。
2025 年初的 o1 仅得 8.81\%，
Claude 3.7 Sonnet 得 8.93\%；
到 2026 年初，
最好的模型提升到约 41\%——
一年内从 9\% 到 41\% 固然令人印象深刻，
但仍有近 60\% 的人类前沿知识题目无法回答。

\textbf{生产代码的鸿沟}。
多项研究发现，LLM 在精心设计的编码基准上得分优异，
但在来自真实 GitHub 仓库的生产代码上
性能\textbf{显著下降}。
合成基准与真实世界代码生成质量之间的相关性
远低于直觉预期——
HumanEval 等基准上的高分
不能可靠地预测模型在复杂工程任务中的表现。

\textbf{基准污染}。
基准数据泄露到训练语料中是一个已知的系统性问题。
多项独立审计发现，
MMLU、GSM8K 和 ARC-Challenge 等广泛使用的基准题目
出现在主要模型的训练数据中。
Llama 4 发布时的基准呈现
引发了社区对结果``粉饰''（cherry-picking）的广泛质疑。

\begin{warnbox}[基准分数 $\neq$ 现实能力]
AI 领域面临一个令人困惑的悖论：
模型在标准评估上的表现持续飙升，
但\textbf{经济影响似乎显著滞后于评估分数的进步}。
基准分数的改善并未按比例转化为
生产力提升或商业价值——
这暗示着当前的评估体系
可能未能捕捉到真正重要的能力维度。
\end{warnbox}

\subsection{这意味着什么？}

这些数据揭示了一个深刻的悖论：
LLM 在\textbf{封闭式、有明确答案}的评估中
已经达到甚至超越博士水平，
但在\textbf{开放式、需要判断力和适应性}的真实世界任务中，
仍然远不及一个有经验的人类专业人士。

更令人不安的是，
我们无法可靠地预测模型会在什么时候、
以什么方式失败。
一个在 GPQA Diamond 上得分 94\% 的模型
可能在一个看似简单但措辞不寻常的问题上
给出完全荒谬的答案。
这种不可预测性的根源
正是本章的核心主题——
\textbf{LLM 的根本性局限}。


%% ============================================================
\section{LLM 的根本性缺陷}
\label{sec:fundamental-flaws}

我们将 LLM 的局限性分为六个维度，
每个维度都指向架构或训练范式层面的根本约束，
而非仅凭更多数据或更大模型就能解决的表面问题。

\subsection{幻觉：统计合理性 $\neq$ 事实准确性}

幻觉（hallucination）是 LLM 最广为人知的缺陷~\cite{alansari2025hallucination}。
模型自信地生成看似合理但与事实不符的内容——
编造不存在的论文引用、虚构法律判例、
杜撰历史事件的细节。

幻觉不是偶然的 bug，而是架构的必然产物：

\textbf{目标函数错位}。
Next-token prediction 优化的是\textbf{统计上的合理性}
而非\textbf{事实上的准确性}。
当模型在 ``某位著名教授发表了关于 X 的开创性论文\_\_\_\_''
这样的上下文中预测下一个词时，
它会生成一个\textbf{听起来像论文标题}的字符串，
而非检索一个\textbf{真实存在}的论文标题。
统计合理性与事实准确性在大多数情况下高度相关——
这就是为什么 LLM 在大部分时间是准确的——
但在两者分离的边缘情况中，
模型总是选择统计合理性。

\textbf{Softmax 瓶颈}。
模型\textbf{必须}在每一步输出一个 token 的概率分布，
无法``弃权''或表达``我不知道''。
即使内部不确定性极高，
softmax 也会强制选出一个最高概率的 token。
这与人类专家的行为形成鲜明对比——
一个诚实的专家会说``我不确定，需要查证''。

\textbf{Exposure bias}。
训练时模型看到的是 ground-truth 前缀，
但推理时必须基于自己生成的 token 继续生成。
一旦早期输出包含轻微错误，
这个错误会\textbf{自回归地级联放大}——
模型基于错误的前缀做出更多预测，
错误像滚雪球一样越滚越大。

\textbf{知识边界的缺失}。
LLM 没有``我知道什么''和``我不知道什么''
的内部表征。
2026 年 3 月来自 METU 的 HypoTermInstruct 研究~\cite{uluoglakci2026hypoterm}
通过让模型学习关于\textbf{假想不存在的术语}的知识，
在 800 次 LoRA SFT 实验中
证明可以训练出``认识论谦逊''——
让模型学会说``我不知道''。
但这种方法尚未在前沿模型上大规模验证。

RAG（检索增强生成）是当前生产环境中
缓解幻觉的主要手段，
但它本质上是一种\textbf{外部补丁}，
并未解决模型内部缺乏事实锚定的根本问题。


\subsection{推理的系统性失败：反转诅咒、组合崩溃与 25 步悬崖}

推理能力是 LLM 最令人兴奋的进步之一——
但系统性的失败模式揭示了
这种推理能力的\textbf{脆弱本质}。

Song、Han 与 Goodman 的综述~\cite{song2026reasoning}
系统分类了 LLM 的推理失败，
将其分为基本型、涌现型和鲁棒性型三类。
以下三个具体的失败模式最具启示性：

\textbf{反转诅咒（Reversal Curse）}~\cite{berglund2023reversal}。
如果模型在训练数据中见过``A 是 B''（例如``Tom Cruise 的母亲是 Mary Lee Pfeiffer''），
它\textbf{无法}回答``B 是 A''（``Mary Lee Pfeiffer 的儿子是谁？''）。
根因直指自回归训练的本质：
模型学习的是\textbf{单向}条件概率 $P(B \mid A)$，
而非对称关联。
这不是一个可以通过更多数据解决的问题——
除非训练数据中显式包含反向表述，
否则模型在架构层面就缺乏逆向推理的能力。

\textbf{Two-Hop 诅咒}~\cite{balesni2025twohop}。
即使微调了 $A \to B$ 和 $B \to C$ 的关系，
LLaMA-3-8B 也\textbf{无法在潜空间中}
可靠地推导 $A \to C$（不使用 CoT 的情况下）。
这说明组合性推理——
将两个独立学到的知识片段组合起来得出新结论——
可能是 Transformer 潜空间中的\textbf{根本性无能力}，
而非偶尔的失败。

\textbf{推理步数的泛化悬崖}~\cite{li2026reasoninghop}。
当推理所需的步数超过训练分布中见过的最大步数时，
模型性能\textbf{骤然崩塌}，
即使底层算法完全相同。
错误集中在特定的 token 位置而非均匀分布，
暗示模型学习的是
``N 步推理模式''的统计模板，
而非可泛化的推理算法。

\begin{whybox}[推理模型（o3/R1）解决了这些问题吗？]
答案是\textbf{部分地}。
CoT（链式推理）可以将 Two-Hop 问题
分解为两个单步推理，
从而绕过潜空间中的组合限制。
但 CoT 本身引入了新的问题：
更长的推理链意味着更多的自回归步骤，
每一步都累积误差。
Gary Marcus 在 2026 年 2 月引用 Song 等人的综述，
认为尽管投入了近万亿美元，
LLM 的推理仍然``深度有缺陷''。
Marcus、Quattrociocchi 与 Capraro
在 \textit{Nature} 回应中更尖锐地指出~\cite{marcus2026nature}：
``Statistical approximation $\neq$ general intelligence''——
高基准分数来自对训练数据的精密模式插值，
而非真正的推理。
\end{whybox}


\subsection{规划的架构性限制}

Wang 等人~\cite{wang2026whyreasoning}
（Notre Dame/Stanford/Edinburgh，2026 年 1 月）
提出了一个尖锐的论点：
\textbf{LLM 的规划失败不是 bug，而是架构限制}。

核心论证如下：
Step-by-step 推理产生的是``逐步贪心策略''——
模型在每一步选择当前看起来最好的动作，
但这种策略在需要考虑\textbf{延迟后果}的场景中失败。
推理（评估当前步骤的正确性）
和规划（优化\textbf{未来}状态序列）
是\textbf{根本不同的认知能力}。

SokoBench 上的实验证实了这一理论~\cite{monti2025sokobench}：
LLM 的性能在任务需要超过约 \textbf{25 步}规划时骤降，
增大模型规模\textbf{不能消除}这一``悬崖''。

Li 等人~\cite{li2026limitedspace}
（2026 年 2 月）从信息论角度提供了补充解释：
推理计算预算存在\textbf{最优区间}——
过度规划（更多 CoT token）反而\textbf{损害}推理能力。
纯过程性推理会``信息停滞''，
模型耗尽了上下文中可提取的新信息。
这解释了为什么简单地增加 thinking budget
并不总能改善结果。


\subsection{灾难性遗忘：可塑性-稳定性困境}

LLM 的权重是\textbf{静态}的——
一旦训练完成，知识就被冻结在参数中。
微调可以注入新知识，
但往往以牺牲已有能力为代价。
这就是\textbf{灾难性遗忘}。

Imanov~\cite{imanov2026catastrophic}
（2026 年 1 月）首次从参数级别进行了机制分析，
揭示了遗忘的直接原因：
新知识在共享参数子空间中\textbf{覆盖}旧知识。
当两个任务依赖相同的参数区域时，
学习一个就必然损害另一个。

Google Research 的 Nested Learning~\cite{behrouz2025nested}
（2025 年 11 月）
将模型视为一组嵌套优化问题，
声称可以``完全避免''灾难性遗忘，
但该方法仍处于早期阶段，
尚未在前沿 LLM 规模上验证。

可塑性（学习新知识的能力）与稳定性（保留旧知识的能力）
之间的张力是\textbf{梯度学习的根本属性}，
而非特定架构的缺陷。
目前没有方案能在前沿模型规模下
令人满意地解决这一困境。


\subsection{鲁棒性的脆弱}

LLM 的表现对输入的微小变化
呈现出令人不安的敏感性。

\textbf{Brittlebench}~\cite{romanou2026brittlebench}
（Meta/EPFL，2026 年 2 月）
对保持语义不变的扰动
（拼写错误、重新措辞、句式变换）
进行了系统测试。
结果触目惊心：
这些微小扰动使模型性能下降高达 12\%，
更关键的是，
\textbf{单个扰动在 63\% 的情况下改变了模型之间的相对排名}。
这意味着排行榜上的模型排名
在很大程度上取决于特定的 prompt 措辞，
而非真正的能力差异。

\textbf{Prompt 注入}
在 OWASP LLM Top 10（2025 v2.0）中排名第一~\cite{owasp2026llm}。
Hu 等人~\cite{hu2026promptinjection}
（2026 年 2 月）证明对抗训练存在根本缺陷——
将提示改为过去时态、翻译成其他语言等
\textbf{简单的分布内变换}就能绕过防御。
根本原因在于 LLM 的核心能力——
\textbf{指令遵循——本身就是漏洞}：
模型无法从架构层面区分
``来自用户的指令''和``来自数据的文本''。


\subsection{理解 vs.\ 模拟：一场远未结束的辩论}

LLM 是否真正``理解''语言和世界？
2026 年的辩论已经高度细化。

\textbf{``仍然是模式匹配''阵营}指出：
Song 等人~\cite{song2026reasoning}系统记录的推理失败
（反转诅咒、组合崩溃、对表面扰动的脆弱性）
表明 LLM 依赖统计模式而非抽象理解。
Marcus 等人~\cite{marcus2026nature}的 \textit{Nature} 回应
更直接地宣称
高基准表现来自``精密的统计近似''。

\textbf{``不止是模式匹配''阵营}也有有力的论据。
Margaret Mitchell——
``随机鹦鹉''（Stochastic Parrots）论文的联合作者——
在 2026 年 3 月撰文指出，
原始框架被误用于描述当前系统。
Anthropic 的机制可解释性研究
发现模型内部存在\textbf{抽象特征}
（如``欺骗''``谄媚''``代码正确性''等概念）~\cite{lindsey2025biology}，
暗示模型的内部表征超越了表面模式匹配。

\textbf{新兴的中间立场}或许最接近真相：
LLM 可能占据了一个\textbf{全新的认知位置}——
具备强大的\textbf{结构性理解}
但缺乏\textbf{经验性锚定}。
它们可以操纵概念之间的关系，
但这些概念并不植根于物理世界的直接经验。
这迫使我们重新思考``理解''本身的定义。


%% ============================================================
\section{黑箱困境：可解释性为何追不上能力增长}
\label{sec:interpretability}

尽管 LLM 展现了令人震惊的能力，
我们对其\textbf{内部工作机制}的理解
仍然远远落后于能力本身的增长。
Dario Amodei 在 2025 年 4 月
将可解释性称为``技术史上前所未有的挑战''，
MIT Technology Review 将其列为 2026 年突破性技术之一。

\subsection{机械可解释性的突破}

2025--2026 年，
机械可解释性（mechanistic interpretability）
取得了几项重要进展。

\textbf{电路追踪（Circuit Tracing）}。
Anthropic 在 2025 年 3 月发布了两篇里程碑论文。
Ameisen 等人~\cite{ameisen2025circuits}
引入了\textbf{归因图（attribution graphs）}，
使用跨层转码器（Cross-Layer Transcoders, CLTs）
替代 MLP 层，
创建了模型计算过程的可解释图形描述。
Lindsey 等人~\cite{lindsey2025biology}
对\textbf{生产模型} Claude 3.5 Haiku 进行了内部机制调查，
发现了多步推理电路、语言内部表示、安全机制等结构。
这是首次在规模化的、
部署中的商用模型上成功追踪内部计算路径。

\textbf{叠加假说（Superposition Hypothesis）}。
Liu、Liu 与 Gore（MIT）~\cite{liu2025superposition}
（NeurIPS 2025 Best Paper Runner-Up）
提出了一个重要的理论结果：
表示叠加是神经网络缩放定律的\textbf{核心驱动力}。
在强叠加状态下，
loss 与模型维度成反比（$L \sim 1/m$），
且与数据分布无关。
这为``为什么更大的模型更好''
提供了第一性原理的解释——
更多的维度意味着更少的叠加干扰，
意味着更精确的特征表示。

\subsection{可解释性的根本障碍}

然而，29 位研究者跨 18 个机构
共同撰写的开放问题论文~\cite{nanda2025openproblems}
（Sharkey、Chughtai 等 29 位研究者，
2025 年发表于 TMLR）清晰地勾勒了
该领域面临的核心挑战：

\textbf{``特征''缺乏严格定义}。
机械可解释性的核心概念——``特征''（feature）——
至今没有数学上的严格定义。
研究者依赖直觉和案例分析
来判断一个被提取的方向是否是有意义的特征，
这使得整个领域缺乏可证伪的标准。

\textbf{计算不可解性}。
许多可解释性查询
被证明是 NP-hard 或更难的计算问题。
当模型规模达到数千亿参数时，
即使近似方法也面临严峻的可扩展性挑战。

\textbf{CoT 不等于透明}。
Anthropic \textbf{自己}的研究（2025 年 3 月）
指出了一个令人不安的发现：
CoT 推理链可能\textbf{不反映模型的实际内部计算}——
模型可能在``为既定结论编造合理论证''。
换句话说，即使推理模型``展示了思考过程''，
这个外化的过程也不一定是真实的内部计算路径。

\textbf{安全性的根本限制}。
Neel Nanda（2025 年 5 月）直言不讳：
``Interpretability will not reliably find deceptive AI''——
即使使用当前最好的可解释性方法，
也\textbf{无法保证}检测到模型中的欺骗行为。
Anthropic 设定的目标是
2027 年前``可靠检测大多数 AI 模型问题''，
但许多研究者认为这\textbf{过于乐观}。

\subsection{理论理解的曙光}

在更基础的层面，
2025--2026 年对 LLM 的理论理解取得了显著进展。

\textbf{Grokking（突然泛化）的理论理解}。
Grokking 现象——模型在训练后期``突然''
从记忆切换到泛化——
正在获得越来越深入的理论解释。
多项研究将其与\textbf{相变}类比~\cite{cullen2026grokking}：
不同的学习策略
（记忆 vs.\ 泛化）
对应参数空间中的不同吸引盆，
Grokking 是两个吸引盆之间的跃迁。
这些工作标志着对 Grokking 的理解
正在从定性描述走向定量分析。

\textbf{LLM 与 Solomonoff 归纳}。
Wan 与 Mei~\cite{wan2025solomonoff}
（投稿 ICLR 2026）
首次将 LLM 与算法信息论正式连接：
LLM 的训练近似 Solomonoff 先验，
next-token prediction 实现近似 Solomonoff 归纳。
这提供了一个深刻的理论视角——
LLM 的泛化能力可能源于
对训练数据的\textbf{最优压缩}，
而非对特定任务的``学习''。
Pan 等人~\cite{pan2025compression}
（NeurIPS 2025 Spotlight）
用 Kolmogorov 结构函数进一步支持了这一观点。

\textbf{相变与信息论}。
\textit{Nature} npj AI（2026 年 2 月）
发表的研究表明，
LLM 在压缩阈值处展现\textbf{相变}——
这意味着模型能力的``涌现''可能
确实是物理学意义上的相变，
而非简单的量变到质变。

\begin{corebox}[理论理解的意义]
这些理论进展的价值
不在于立即改善模型性能，
而在于提供\textbf{可预测性}。
如果我们能从第一性原理理解
模型为什么在某个规模上获得某种能力，
我们就能更可靠地预测
下一代模型将展现什么能力——
和什么风险。
\end{corebox}


%% ============================================================
\section{因果推断：大语言模型缺失的一环}
\label{sec:causality}

在 LLM 的所有局限性中，
\textbf{因果推理的缺失}
可能是最根本、最深远的。
它不仅解释了
为什么 LLM 在某些看似简单的问题上失败，
更指向了当前范式的\textbf{理论天花板}。

\subsection{Pearl 的因果阶梯与 LLM 的位置}

Judea Pearl 的\textbf{因果阶梯}
（Ladder of Causation）~\cite{pearl2009causality}
定义了三个层级的因果推理：

\textbf{第一层：关联（Association）}——``X 和什么共同出现？''
纯观测性的统计相关。
$P(Y \mid X)$。

\textbf{第二层：干预（Intervention）}——``如果我做 X 会怎样？''
需要 do-calculus：$P(Y \mid \text{do}(X))$。
这与条件概率\textbf{根本不同}——
$P(\text{地湿} \mid \text{下雨})$ 是关联，
$P(\text{地湿} \mid \text{do(洒水)})$ 是干预。
知道下雨与地湿的相关性
不能告诉你洒水是否会让地湿。

\textbf{第三层：反事实（Counterfactual）}——
``如果 X 不同，会发生什么？''
需要推理替代世界：
``如果那天没下雨，地会湿吗？''

Pearl 的论断是尖锐的~\cite{pearl2025aaas}：
LLM 在数学上被限制在第一层。
它们训练于观测文本数据的 next-token prediction，
学习的是 $P(\text{token} \mid \text{context})$，
这是\textbf{纯关联分布}。
更多的参数和数据不能
将系统从第一层升级到第二层。
LLM 看似正确的因果回答
是在``检索''训练文本中\textbf{人类写入的因果知识}，
而非自行执行因果推理。

\subsection{证据与反证}

然而，领域对此存在真正的分歧。

\textbf{支持 Pearl 的证据}。
当表面语义线索被中和后
（例如将变量名替换为无意义的符号），
即使最强的 2025 年模型
在结构相同的因果任务上
\textbf{性能骤降至接近随机水平}。
多项研究发现，
当表面语义线索被控制或移除后，
LLM 在因果推理任务上的表现\textbf{大幅下降}，
在某些形式化因果基准上
仅略高于随机基线。
这暗示高准确率
可能部分源于语义模式记忆
而非真正的因果推理。

\textbf{反对 Pearl 的证据}。
Kiciman 等人~\cite{kiciman2024causal}
（Microsoft Research）
发现 GPT-4 在多个因果推理基准上
显著超越传统方法——
在成对因果发现等任务上
达到了令人惊讶的高准确率。
他们认为 LLM 已经以\textbf{功能性方式}
内化了因果知识。

如何调和这些矛盾？
一个可能的解释是：
LLM 在因果推理上表现出\textbf{两层能力}——
当任务的因果结构与训练数据中常见的模式匹配时，
模型可以正确``检索''并应用因果知识（高分）；
当任务的因果结构\textbf{新颖}
或与常见模式冲突时，
模型退化为统计猜测（低分）。
这与``精密模式插值''的假说一致——
模型擅长在训练分布\textbf{内}的因果推理，
但无法泛化到分布\textbf{外}的新颖因果结构。

\subsection{因果推理为何未被大规模应用}

截至 2026 年初，
\textbf{没有任何已知的前沿 LLM
在预训练目标中整合了显式的因果推理机制}。
原因是多方面的：

\textbf{计算不可行性}。
结构因果模型（SCM）需要有向无环图（DAG），
变量数的超指数增长使得
在亿级隐变量规模下构建 DAG 不可能。
前沿 LLM 的参数空间远超
任何因果发现算法的处理能力。

\textbf{缺乏监督信号}。
训练语料中因果关系\textbf{未被标注}。
文本中写着``抽烟导致肺癌''，
但这个``导致''对模型来说
只是一个高频出现在``抽烟''和``肺癌''之间的 token，
而非一个显式的因果边。
干预数据（controlled experiment data）极其稀缺。

\textbf{训练目标不兼容}。
Next-token prediction 在本质上是关联性的
$P(Y \mid X)$。
引入 do-calculus 需要\textbf{根本不同的目标函数}——
一个能够区分
``X 和 Y 共同出现''与``X 导致 Y''的目标。
目前不存在这样的目标函数
能够在万亿 token 规模上高效训练。

\textbf{``足够好''论点}。
当前 LLM 已经从文本中吸收了大量人类书写的因果知识，
在大多数实际应用中表现``足够好''。
商业压力使得实验室更倾向于
优化已有范式而非探索根本替代方案。

\textbf{可识别性挑战}。
仅从观测数据识别因果结构
需要强假设（如 faithfulness、causal sufficiency），
这些假设在规模化场景中\textbf{难以验证}。

\subsection{混合架构的曙光}

尽管因果推理未进入预训练，
后训练和混合架构方面出现了令人鼓舞的进展。

\textbf{公理化训练}。
Vashishtha 等人~\cite{vashishtha2025axiomatic}
（ICML 2025）
证明了一个关键结果：
67M 参数的 Transformer
从头训练因果公理（传递性、d-分离等）的示例后，
\textbf{能够泛化到未见过的更复杂图结构}。
当扩展到 LLaMA-3-8B-Instruct 微调时，
因果 benchmark 性能显著提升。
这说明 Transformer \textbf{架构本身}
并不排斥因果推理——
问题在于\textbf{训练数据和目标}，
而非模型架构。

\textbf{因果发现可以通过 NTP 学习}。
Butkus 与 Kriegeskorte~\cite{butkus2026causaldiscovery}
（Columbia，NeurIPS 2025）
做出了一个令人惊讶的发现：
GPT 风格的 Transformer
在\textbf{干预数据}上训练后，
可以\textbf{同时发现因果结构并回答反事实查询}。
从残差流（residual stream）中
可以解码出隐式的``心理'' SCM。
\textbf{关键含义}：
next-token prediction 目标与因果推理
\textbf{并非本质不兼容}——
前提是训练数据包含干预信息。

\textbf{因果奖励建模}。
\textbf{因果奖励建模}方面的研究~\cite{crome2025deepmind}
使用反事实数据增强
来区分 RLHF 中的真实质量驱动因素与虚假相关
（如回复长度与质量的混淆）。
早期实验显示，
因果感知的奖励模型在安全性和推理任务上
均有显著提升。
这代表了将因果推理应用于
LLM 对齐流水线的新方向。

\textbf{Language Agents Meet Causality}。
Gkountouras 等人~\cite{gkountouras2025languageagents}
（ICLR 2025 Poster）
将因果表示学习（CRL）与 LLM Agent 结合：
CRL 模块学习环境的因果结构，
LLM 用此结构进行因果感知的规划。
这种混合架构绕过了 Pearl 指出的根本限制——
不要求 LLM \textbf{自身}执行因果推理，
而是为它提供一个\textbf{外部因果引擎}。

\textbf{DEMOCRITUS}。
Mahadevan 等人~\cite{mahadevan2025democritus}
（Adobe/UMass，2025 年 12 月）
提出了一种反向思路：
不是让 LLM 学习因果推理，
而是\textbf{从 LLM 中提取因果知识}
来构建大规模因果图。
DEMOCRITUS 系统通过迭代自校正
从 LLM 文本输出中提取因果关系，
使用 ``Geometric Transformer'' 架构
组装成结构化的因果模型。

\subsection{Sch\"{o}lkopf 的因果表示学习}

Bernhard Sch\"{o}lkopf（MPI-IS T\"{u}bingen / ELLIS Institute）
是因果性与现代机器学习交叉领域
最具影响力的学者之一。

他的核心思想是\textbf{独立因果机制（ICM）原则}~\cite{scholkopf2021causalrep}：
世界的生成过程由\textbf{独立自主的模块}组成，
改变一个机制不会影响其他机制。
这个原则为\textbf{因果表示学习（CRL）}
提供了理论基础——
如果我们能从数据中学习到
满足 ICM 原则的表示，
那么这些表示就具有真正的因果意义，
而非仅仅是统计相关。

2024--2025 年的进展
正在将这一理论愿景推向实践：
从因果到概念的表示学习~\cite{rajendran2024causalconcept}
（NeurIPS 2024，与 Buchholz、Aragam 等人合作）
桥接了 CRL 和基础模型中的概念学习。
CLadder benchmark~\cite{jin2023cladder}
（NeurIPS 2023）
为 LLM 的因果推理评估建立了标准。

\begin{keybox}[因果推理的未来]
2026 年初的领域共识正在形成：
\textbf{最有前途的路径不是
让 LLM 自身成为因果推理器，
而是将因果模型作为外部组件
与 LLM 结合}。
这与 Agent 范式天然契合——
LLM 负责语言理解和生成，
因果引擎负责干预推理和反事实分析，
二者通过工具调用接口连接。
这可能是超越 Pearl 所指出的
关联-干预鸿沟的最实际路径。
\end{keybox}


%% ============================================================
\section{世界模型：理解物理因果的替代路线}
\label{sec:world-models}

第~\ref{ch:outlook}章
已经从视频生成的角度介绍了世界模型。
本节从更根本的技术层面
讨论 Yann LeCun 的替代路线——
为什么他认为自回归 token 预测
\textbf{原则上}无法达到真正的智能？
他提出的 JEPA 架构有什么独特之处？
以及这条路线面临什么样的批评？

\subsection{LeCun 对自回归范式的六大技术批评}

LeCun 对当前 LLM 范式的批评~\cite{lecun2022path}
不是笼统的哲学反对，
而是具体的、可检验的技术论点：

\textbf{批评一：指数误差累积}。
自回归生成中，
每一步的正确率为 $1-\epsilon$，
$N$ 步后的序列正确率为 $(1-\epsilon)^N$。
这意味着即使每一步 99\% 正确，
100 步后正确率也只有 36.6\%。
推理模型的长链推理
从根本上受限于这一指数衰减。

\textbf{批评二：信息带宽差距}。
视觉输入的信息速率约 20 MB/s，
而语言的信息速率不到 12 bytes/s——
差距达到 \textbf{160 万倍}。
一个仅在语言上训练的系统
\textbf{错过了世界信息的 99.9999\%}。
人类的大部分世界知识
来自视觉和触觉经验，
而非语言描述。

\textbf{批评三：离散 token 无法表征连续物理世界}。
物理世界是连续的——
一个球的轨迹、一杯水的流动、一根绳子的弯曲
都是连续变量的函数。
将世界离散化为 token
丢失了连续性中蕴含的物理约束信息。

\textbf{批评四：缺乏因果/物理基础}。
LLM 学到的是统计相关而非因果关系。
``苹果从树上掉下来''在训练数据中频繁出现，
但模型不``理解''这是因为重力——
如果你问它``在月球上苹果会怎样''，
它会基于训练数据中的月球描述回答，
而非从万有引力定律推导。

\textbf{批评五：无持久世界模型}。
LLM 的权重在训练后冻结，
没有一个随交互持续更新的世界状态表示。
每次对话都从空白开始——
模型没有``记忆''昨天桌子上有什么东西。

\textbf{批评六：每 token 固定计算量}。
标准 Transformer 对每个 token
分配\textbf{完全相同}的计算量，
无论是一个简单的冠词``the''
还是一个需要深度推理的数学表达式。
这从根本上限制了模型的推理能力——
复杂问题需要更多的计算，
但架构不允许这种动态分配。

\subsection{JEPA：在表征空间中预测}

LeCun 提出的替代方案是
\textbf{联合嵌入预测架构}
（Joint Embedding Predictive Architecture, JEPA）~\cite{lecun2022path}。

JEPA 的核心创新在一句话中：
\textbf{在学到的抽象表征空间中预测未来状态，
而非在像素或 token 空间中}。

架构由三个组件构成：
（1）\textbf{Context Encoder}——
将当前观测编码为潜在表示；
（2）\textbf{Target Encoder}——
将目标（未来状态）编码为潜在表示，
使用 EMA（指数移动平均）权重更新，
\textbf{不参与梯度训练}；
（3）\textbf{Predictor}——
在潜在空间中预测 Target Encoder 的输出。

与自回归模型的根本区别：
预测\textbf{连续的抽象表征}而非离散 token。
与生成模型（如扩散模型）的区别：
不需要重建每个像素细节——
未来一片树叶向左飘还是向右飘
在像素层面是巨大的不确定性，
但在\textbf{语义层面}是微不足道的
（树叶会落地，不会飞上天）。
JEPA 在语义层面预测，
自然地过滤掉了不相关的像素级变化。

从能量模型（EBM）的角度看，
JEPA 是 EBM 的特例：
能量函数 = 嵌入空间中的预测误差。
低能量表示合理的结果，
高能量表示不合理的结果。
规划就是在能量景观上的\textbf{优化问题}——
找到一系列动作使得
未来状态的能量最小化。

\subsection{V-JEPA 2：视频世界模型的突破}

V-JEPA 2~\cite{vjepa2_2025}（2025 年）
是 JEPA 路线迄今最大规模的验证。
作为视频世界模型，
它采用两阶段训练：
（1）大规模视频和图片的自监督预训练；
（2）少量机器人数据的微调。

关键结果包括：
在视频理解基准上达到 SOTA；
零样本机器人规划在未见过的物体和新环境中
展现了可观的成功率；
推理规划速度显著优于同类方法。

Meta FAIR 同时发布了三个新的物理推理基准，
揭示了当前 AI 距人类物理理解的\textbf{巨大差距}——
即使在 V-JEPA 2 最擅长的领域，
人类仍然大幅领先。

后续的 VL-JEPA 工作
将 JEPA 扩展到视觉-语言领域，
在参数效率和解码速度上
相对标准 VLM 展现了显著优势。
这表明 JEPA 的表征效率
在多模态场景中具有潜力。


\subsection{反驳与辩论：Bitter Lesson 是否适用？}

LeCun 的世界模型路线面临几个有力的反驳：

\textbf{预测屡被证伪}。
LeCun 曾预测``GPT-5000 也无法理解
桌子推动书本移动''，
但现代多模态模型已经可以轻松做到这一点。
每次他设定一个``LLM 永远无法做到''的门槛，
scale up 后的模型就跨过了它。

\textbf{Bitter Lesson}。
Richard Sutton 的``苦涩教训''
指出历史上特定领域的巧妙方法
最终总是被通用的规模扩展击败。
JEPA 嵌入了强烈的\textbf{架构先验}
（在表征空间预测、能量模型框架），
而自回归 LLM 的简洁目标
更符合 Bitter Lesson 的精神。
然而，值得注意的是，
Sutton \textbf{本人}在 2025 年也表示
LLM 不是 Bitter Lesson 的体现——
LLM 的成功可能更多来自数据和工程，
而非算法的通用性。

\textbf{指数误差论可能过于悲观}。
LeCun 的 $(1-\epsilon)^N$ 分析
假设每个 token 的错误概率\textbf{独立且均匀}，
但实际的自回归生成中，
大部分 token（冠词、介词等）的预测几乎无误，
只有少数``关键 token''（推理步骤、事实断言）
对序列正确性有实质影响。
如果仅少数关键 token 承载误差风险，
实际的衰减可能比指数函数温和得多。

\textbf{JEPA 尚未产生变革性应用}。
截至 2026 年初，
JEPA 在视频理解和机器人规划上展现了潜力，
但尚未在\textbf{任何}领域
达到自回归 LLM 的实际影响力。
V-JEPA 2 的零样本机器人规划令人印象深刻，
但距离生产级部署还有相当距离。

\subsection{从学术宣言到产业实践}

世界模型路线正在从学术论文中的理论构想
进入有实质资本支撑的产业化阶段。
Meta FAIR 持续投入 JEPA 系列研究，
V-JEPA 2 和 VL-JEPA 的成果
证明了这条路线在工程上的可行性。
更广泛地看，
投资界对``超越自回归''的替代路线
表现出越来越浓厚的兴趣——
多家初创公司正在探索
世界模型、因果推理和具身智能的商业化。

LeCun 的赌注很明确：
如果自回归 LLM 的局限性
确实是架构性的而非规模性的，
那么率先在世界模型上取得突破的团队
将赢得 AI 的下一个时代。
这是否会重复
``专家系统 vs.\ 神经网络''历史中
LeCun 最终获胜的剧本？
还是会像许多理论上优雅
但实践中不敌暴力规模扩展的方法一样被历史遗忘？
答案可能需要 3--5 年才能揭晓。


%% ============================================================
\section{超越 Transformer：新架构与新范式}
\label{sec:beyond-transformer}

如果当前 LLM 的局限性
确实部分源于 Transformer 架构本身，
那么超越 Transformer 的架构创新
就不仅是学术兴趣，
更是实际需求。
2025--2026 年，
几个方向的进展值得关注。

\subsection{状态空间模型：亚二次计算的回归}

Transformer 的自注意力机制
具有 $O(n^2)$ 的计算复杂度，
这限制了处理极长序列的能力。
状态空间模型（State Space Models, SSM）
提供了一种 $O(n)$ 或亚二次的替代方案。

\textbf{Mamba-3}~\cite{lahoti2026mamba3}
（CMU/Princeton，ICLR 2026 Oral）
是``推理优先''的 SSM 架构，
改进了离散化方法、引入复数动态和 MIMO 更新。
它实现了亚二次计算和\textbf{恒定内存}的推理——
无论输入多长，每生成一个 token 的内存消耗不变。
这与 Transformer 的 KV 缓存
随序列长度线性增长形成鲜明对比。

\textbf{RWKV-7 ``Goose''}~\cite{peng2025rwkv7}
（EleutherAI，COLM 2025）
实现了恒定内存和恒定推理时间/token，
2.9B 参数在多语言任务上达到 3B 级别的 SOTA。
RWKV-7 使用\textbf{广义 delta 规则}，
在标准复杂性假设下
可以执行状态追踪并识别\textbf{所有正则语言}——
这在理论表达能力上超越了
Transformer 的 $\text{TC}^0$ 限制。
完全开源。

\textbf{混合架构}是近期的实用赢家。
AI21 Labs 的 Jamba
将 Transformer、Mamba 和 MoE \textbf{混合}在一起，
398B 参数（94B 活跃），256K 上下文。
这种``取各家之长''的工程哲学
可能比追求单一架构的理论纯粹性更务实——
在需要精确回忆的层使用注意力，
在处理长距离依赖的层使用 SSM。


\subsection{扩散语言模型：自回归的替代}

如果 next-token prediction 的自回归本质
是某些局限性的根源（如反转诅咒、指数误差累积），
那么\textbf{非自回归}的文本生成方法就值得探索。

\textbf{掩码离散扩散模型（MDLM）}~\cite{sahoo2025mdlm}
（Cornell，NeurIPS 2025）
证明了掩码离散扩散
比此前认为的\textbf{更强}——
通过精心设计的噪声调度和参数化，
MDLM 在多个文本生成基准上
接近自回归模型的质量。

\textbf{离散流匹配（Discrete Flow Matching）}~\cite{gat2025dfm}
（Meta FAIR，NeurIPS 2025）
为离散数据提供了通用的流匹配范式。
Apple 的后续工作 FS-DFM（ICLR 2026）
在仅需\textbf{少量步数}的条件下
实现了接近自回归质量的\textbf{并行生成}。

这些方法的潜在优势在于：
扩散/流匹配模型天然具有
\textbf{全局一致性}——
它们在整个序列上
同时进行去噪/细化，
而非像自回归那样
从左到右逐步``承诺''每一个 token。
这可能缓解反转诅咒等单向生成导致的问题。

\textbf{能量模型视角}。
Stanford/NVIDIA 的能量基扩散语言模型~\cite{xu2025energydiffusion}
（ICLR 2025）
和 NRGPT~\cite{dehmamy2026nrgpt}（ICLR 2026 Poster）
将 GPT 与能量模型统一，
推理被建模为能量景观上的\textbf{梯度下降}——
与 LeCun 的 JEPA 在哲学上异曲同工。


\subsection{测试时训练与连续思维}

\textbf{测试时训练（Test-Time Training, TTT）}
代表了另一种根本性的创新：
不是在训练后冻结模型，
而是在\textbf{推理时继续学习}。

Tandon 等人~\cite{tandon2025ttt}
（Stanford/Berkeley/NVIDIA，2025 年 12 月）
将长上下文建模重新定义为\textbf{持续学习问题}。
标准 Transformer + 滑动窗口注意力，
但在推理时通过 next-token prediction \textbf{继续训练}，
将上下文信息\textbf{压缩到权重中}。
这从根本上改变了
``上下文窗口''的含义——
不再是一个固定大小的缓冲区，
而是一个可以持续吸收信息的\textbf{学习系统}。

Aky\"{u}rek 等人~\cite{akyurek2025ttt}
（MIT）
展示了 TTT 在少样本学习中的惊人效果：
ARC 基准上准确率\textbf{大幅提升}，
接近人类平均水平。

\textbf{连续思维机器（Continuous Thought Machines）}~\cite{sakana2025ctm}
（Sakana AI，NeurIPS 2025 Spotlight）
重新引入了\textbf{神经元级时间处理}
和\textbf{同步}作为核心表示机制。
与标准 Transformer 中
所有神经元同步更新不同，
CTM 允许不同神经元以不同速率演化，
类似于生物大脑中的时间编码。
在迷宫求解、图像分类等任务上，
CTM 展现了\textbf{涌现式的逐步推理}——
模型自发地学会了
``先观察全局，再聚焦细节''的策略，
而这种策略从未被显式训练。

\textbf{TopoNets}~\cite{rajkumar2025toponets}
（Georgia Tech，ICLR 2025 Spotlight）
模仿大脑的拓扑组织，
将神经元排列在 2D 皮层结构上，
效率提升 20\%+ 几乎无性能损失。
\textit{Nature Communications}（2025）
上发表的工作进一步展示了
模拟大脑区域功能分工的
\textbf{模块化 Agent 架构}。


%% ============================================================
\section{从商业狂热回到学术范式}
\label{sec:back-to-academia}

2024--2025 年是 LLM 商业化的``淘金热''。
数十亿美元涌入更大的模型、更多的 GPU 集群、
更快的推理服务。
``规模就是一切''是这一时期的主旋律。

但 2025 年下半年开始，
多个信号表明这一范式正在遭遇瓶颈。

\textbf{预训练规模扩展的减速}。
正如第~\ref{ch:outlook}章所讨论的，
高质量文本数据正在接近耗尽，
算力-性能边际递减使得
下一个 10\% 的改善需要数百亿美元。
业界越来越多的声音承认，
纯粹的规模扩展可能不足以实现 AGI，
\textbf{算法层面的根本突破}不可或缺。

\textbf{算法创新的回归}。
回顾 2025--2026 年最有影响力的进步，
许多并非来自更大的模型或更多的 GPU：
\textbf{真正的飞跃来自算法创新}。
DeepSeek-R1 的 GRPO 不需要更多的 GPU，
却让模型学会了自我反思和推理验证。
这种``少投入、大产出''的算法创新
与``更多 GPU、更多数据''的暴力扩展
代表了两种根本不同的进步路径。

\textbf{Sutskever 的``时代转变''论}。
Ilya Sutskever 明确宣称
AI 领域正从``缩放时代''进入``研究时代''——
预训练数据墙正在逼近，
未来的突破将更多来自
架构创新、训练方法创新和理论理解，
而非简单地堆叠更多计算。

这一转变对学术界是一个利好信号。
过去几年，
``只有拥有 10 万张 GPU 的大公司
才能推动 AI 前沿''
几乎成为共识。
但如果未来的关键进步
来自算法创新而非规模堆叠，
那么\textbf{一个有洞察力的小团队}
可能比一个有无限预算的大公司
做出更大的贡献。
DeepSeek 已经证明了这一点，
AMI Labs 正在下这个赌注。

\begin{corebox}[学术界可以贡献什么？]
回顾本章讨论的前沿方向：
\textbf{因果推理}的理论基础完全来自学术界
（Pearl、Sch\"{o}lkopf）；
\textbf{机械可解释性}的核心突破来自 Anthropic
但理论框架来自 MIT、Oxford 等大学；
\textbf{状态空间模型}起源于 CMU 和 Stanford 的博士论文；
\textbf{Grokking 和相变}的理论理解
来自纯粹的学术好奇心驱动的研究。
在``规模为王''的时代，
这些研究看似``不实用''。
但如果规模扩展真的接近瓶颈，
它们可能成为
打破瓶颈的\textbf{关键钥匙}。
\end{corebox}


%% ============================================================
\section{本章总结}
\label{sec:limitations-summary}

让我们回顾本章的核心论点。

\textbf{LLM 同时是
人类创造的最强大的智能工具
和最不透明的技术系统之一}。
它们在封闭式基准上超越博士，
在开放式真实任务中
失败率高达 76--82\%。
这个悖论不是
可以通过更多数据或更大模型解决的表面问题——
它指向了
自回归 token 预测范式的\textbf{深层架构限制}。

幻觉源于目标函数优化统计合理性而非事实准确性。
推理失败源于单向条件概率无法支持对称和组合推理。
规划受限于逐步贪心策略无法处理长程后果。
因果推理的缺失
源于关联分布无法表达干预和反事实。
这些不是偶然的 bug，
而是\textbf{数学上可以论证的架构约束}。

同时，学术前沿正在探索多条可能的突破路径：
因果表示学习为 LLM 提供外部因果引擎；
JEPA 在表征空间中预测以避免像素级的不确定性；
状态空间模型提供亚二次复杂度的长序列处理；
扩散语言模型挑战自回归生成的垄断地位；
测试时训练模糊了训练和推理的边界；
机械可解释性逐步揭开黑箱的面纱。

\textbf{从``缩放时代''到``研究时代''的转变}
可能是 2025--2026 年最重要的范式变化。
当规模扩展接近物理和经济极限时，
算法创新、架构创新和理论理解
将重新成为推动 AI 前沿的主要力量。
学术界——
而非仅仅是拥有十万张 GPU 的大公司——
可能再次成为突破的发源地。

如果用一句话总结本章的核心信息：
\textbf{LLM 的能力令人敬畏，
但我们不能因为敬畏而停止追问——
它们为什么工作？
它们何时会失败？
以及，有没有更好的方式？}
